\section{What We Have Right Now}\label{sec:current}

\begin{plainwords}
The Fanela Planner v11 is one HTML file (about 300 KB) that runs inside a web browser.
It has been improved at least eleven times, and each version added real factory rules:
how orders are recorded, how sizes are allocated, how artwork is approved, how embroidery
samples block production, how shipments are tracked.

It works -- but only inside the browser that opened it. Data lives in the browser's local
storage. Accounts and passwords are stored locally. Attachments are limited to 2 MB.
There is no sync between computers, no central backup, and no server to stop a user from
bending the rules. Clearing browser storage deletes everything.
\end{plainwords}

\subsection*{What the prototype does well}

\begin{itemize}
  \item Captures years of hard-won operational detail that would take months of workshops
        to rediscover.
  \item Encodes complex controls: customer snapshots, artwork approval, embroidery swatch
        gating, department queues, dispatch records, append-only stock history.
  \item Honest about its limits -- for example it states clearly that DPD booking is
        manual and there is no live API.
\end{itemize}

\subsection*{What limits it}

\begin{tabularx}{\textwidth}{@{}l X@{}}
\toprule
Limitation & Consequence \\
\midrule
Browser-local storage & One machine, one browser. No teamwork, no sync \\
Local accounts and passwords & No real identity, no multi-factor, no central revocation \\
Client-side permissions & Anyone with developer tools can bypass role restrictions \\
Manual backups & Easy to lose everything; no point-in-time recovery \\
Attachments in browser (2 MB) & Photos and labels barely fit; lost with the browser \\
No concurrency handling & Two users editing the same job silently overwrite each other \\
Audit stored locally & Not tamper-resistant; can vanish with the browser \\
DPD and Xero & Manual, disconnected processes \\
\bottomrule
\end{tabularx}

\begin{warningbox}
Hosting the existing HTML file on a web server would solve none of the above. It would
still be browser-local storage with no multi-user safety. That path is explicitly ruled
out (see anti-patterns, Section~\ref{sec:migration}).
\end{warningbox}
